%========================================
% 9 | Saṅghābhigīti
%========================================
% title-en: Chant in Praise of the Sangha
\section*{e9 | Saṅghābhigīti}

\begin{chantsubsection}
  \chantnote
    {\{Handa mayaṃ saṅghābhigītiṃ karomase\}}
  \chantnotetrans
    {\{Now let us recite a chant in praise of the Saṅgha.\}}
\end{chantsubsection}

\chantgap

% ---- Opening praise ----
\begin{chantsubsection}
  \chantline
    {Saddhammajo supaṭipattiguṇādiyutto}
    {Born of the True Dhamma, endowed with virtues such as good practice,}

  \chantline
    {Yoṭṭhabbidho ariyapuggalasaṅghaseṭṭho}
    {the noble Saṅgha, supreme among the eight kinds of noble persons,}

  \chantline
    {Sīlādidhammapavarāsayakāyacitto}
    {with body and mind established in the excellent abode of virtue and other qualities,}

  \chantline
    {Vandāmahaṃ tamariyānagaṇaṃ susuddhaṃ}
    {I revere that pure and noble assembly.}
\end{chantsubsection}

\newpage

% ---- Main praise (two pādas per line) ----
\begin{chantsubsection}
  \chantlinepair
    {Saṅgho yo sabbapāṇīnaṃ}
    {Saraṇaṃ khemamuttamaṃ}
    {The Saṅgha is for all beings the supreme, secure refuge.}

  \chantlinepair
    {Tatiyānussatiṭṭhānaṃ}
    {Vandāmi taṃ sirenahaṃ}
    {As the third foundation of recollection, I bow to it with my head.}

  % footnote 1
  \chantlinepair
    {Saṅghassāhasmi dāso\footnotemark{} va}
    {Saṅgho me sāmikissaro}
    {I am a servant of the Saṅgha; the Saṅgha is my lord and master.}
  \footnotetext{For ladies: change \textit{dāso} to \textit{dāsī}.}

  \chantlinepair
    {Saṅgho dukkhassa ghātā ca}
    {Vidhātā ca hitassa me}
    {The Saṅgha destroys suffering and brings about my welfare.}

  \chantlinepair
    {Saṅghassāhaṃ niyyādemi}
    {Sarīrañjīvitañcidaṃ}
    {To the Saṅgha I dedicate this body and this very life.}

  % footnote 2
  \chantlinepair
    {Vandantohaṃ\footnotemark{} carissāmi}
    {Saṅghassopaṭipannataṃ}
    {Paying homage, I will live in accordance with the Saṅgha’s good practice.}
  \footnotetext{For ladies: change \textit{Vandantohaṃ} to \textit{Vandantīhaṃ}.}

  \chantlinepair
    {Natthi me saraṇaṃ aññaṃ}
    {Saṅgho me saraṇaṃ varaṃ}
    {I have no other refuge; the Saṅgha is my excellent refuge.}

  \chantlinepair
    {Etena saccavajjena}
    {Vaḍḍheyyaṃ satthu sāsane}
    {By this truthful utterance, may I grow in the Teacher’s dispensation.}

  % footnote 3 (note: first pāda stands alone in the source)
  \chantlinepair
    {Saṅghaṃ me vandamānena\footnotemark{}}
    {Yaṃ puññaṃ pasutaṃ idha}
    {By my paying homage to the Saṅgha, whatever merit is produced here,}
  \footnotetext{For ladies: change \textit{vandamānena} to \textit{vandamānāya}.}

  \chantlinepair
    {Sabbepi antarāyā me}
    {Māhesuṃ tassa tejasā}
    {may all my obstacles be dispelled through its power.}
\end{chantsubsection}

% ---- Gesture note ----
\chantgap
\begin{chantsubsection}
  \chantnote
    {\textit{--- bow, chanting softly ---}}

\chantgap

% ---- Confession verse ----
  \chantline
    {Kāyena vācāya va cetasā vā}
    {By body, speech, or mind,}

  \chantline
    {Saṅghe kukammaṃ pakataṃ mayā yaṃ}
    {whatever wrong I have done toward the Saṅgha,}

  \chantline
    {Saṅgho paṭiggaṇhatu accayantaṃ}
    {may the Saṅgha accept my confession,}

  \chantline
    {Kālantare saṃvarituṃ va saṅghe}
    {so that in the future I may restrain myself with regard to the Saṅgha.}
\end{chantsubsection}
